
O mapeamento de texturas (\textit{texture mapping}) é uma técnica fundamental da Computação Gráfica para associar informações bidimensionais à superfície de objetos tridimensionais. Na sua forma tradicional, as coordenadas \(U\) e \(V\) estabelecem a correspondência entre pontos de uma superfície 3D e posições em uma textura 2D. Heckbert \cite[p.~1]{heckbert1986} apresenta a técnica como uma forma de aumentar a riqueza visual das imagens ``com um aumento relativamente pequeno de computação'', sem exigir que todos os detalhes sejam modelados geometricamente.

A partir dessa base, o mapeamento pode ser modificado durante a execução. No \textit{UV mapping dinâmico}, as coordenadas sofrem translação, escala, rotação ou deformações que dependem do tempo, de modo que a textura se move sobre uma malha parada. No \textit{mapeamento procedural}, a cor ou a própria coordenada é calculada por uma função matemática, em vez de lida de uma imagem \cite[p.~1]{dong2020}.

Este trabalho apresenta uma revisão bibliográfica da técnica, apoiada em três levantamentos: o trabalho clássico de Heckbert sobre texture mapping \cite{heckbert1986}, o levantamento de Dong et al.\ sobre geração procedural de texturas bidimensionais \cite{dong2020} e o de Smelik et al.\ sobre modelagem procedural de mundos virtuais \cite{smelik2014}. A revisão é acompanhada de uma implementação em Vulkan (pasta \texttt{código/}), que renderiza um modelo glTF em cinco modos: textura original, UV dinâmico, padrão procedural, combinação dos dois e grade de depuração do UV. A implementação é usada ao longo do texto como exemplo concreto de cada equação, de cada recurso do Vulkan e de cada limitação discutida.

O documento está organizado da seguinte forma:
\begin{itemize}
    \item o Capítulo~\ref{chap:fundamentacao} apresenta os fundamentos e discute quando usar cada abordagem, suas limitações e seu custo computacional;
    \item o Capítulo~\ref{chap:equacoes} apresenta as equações na ordem em que são executadas no pipeline, indicando o estágio e a frequência de cada uma;
    \item o Capítulo~\ref{chap:vulkan} trata da relação com o Vulkan e da arquitetura do pipeline: diagrama, recursos, descriptor sets, push constants, render pass, sincronização e otimizações;
    \item o Capítulo~\ref{chap:shaders} apresenta os shaders comentados, do pseudocódigo à versão final;
    \item o Capítulo~\ref{chap:aplicacoes} analisa sete produtos reais que usam a técnica;
    \item o Capítulo~\ref{chap:jogos} analisa \textit{Rocket League} e \textit{Marvel's Spider-Man Remastered}, com capturas próprias marcadas, vídeos do grupo e uma demonstração na implementação;
    \item o Capítulo~\ref{chap:resumos} resume as três publicações, e o Capítulo~\ref{chap:conclusao} conclui o trabalho.
\end{itemize}
